\section{Apéndice B: Scripts de Modelado y Simulación en Python}
\label{sec:apendice_simulacion}

En este apéndice se presentan los scripts de simulación y deducción cuantitativa desarrollados en Python 3 para validar analítica y computacionalmente el enlace inductivo a \SI{125}{\kilo\hertz}. Se emplean las bibliotecas científicas \texttt{numpy}, \texttt{scipy} y \texttt{matplotlib}.

\subsection{Deducción Cuantitativa de Parámetros de Enlace}

El script \texttt{calcular\_parametros\_enlace.py} implementa las deducciones formales de ganancia pasiva cargada, ganancia activa considerando el polo de corte por producto ganancia-ancho de banda ($\text{GBW} = \SI{1.0}{\mega\hertz}$), balance de potencia activa y reactiva, y presupuesto temporal de modulación.

\begin{lstlisting}[style=custompy, caption={Cálculo analítico de parámetros de ganancia, potencia e ISI (\texttt{informe\_ieee/datos/calcular\_parametros\_enlace.py})}]
#!/usr/bin/env python3
"""
Validacion Cuantitativa para Enlace Inductivo a 125 kHz
- Ganancia pasiva y activa (con correccion por GBW y carga)
- Balance de potencia en transmision (TX) y recepcion (RX)
- Analisis temporal, interferencia intersimbolica (ISI) y throughput
"""

import math

def main():
    print("=" * 70)
    print("  VALIDACION CUANTITATIVA DEL ENLACE INDUCTIVO A 125 kHz")
    print("=" * 70)

    # 1. Parametros base del sistema
    f0 = 125.04e3          # Hz
    w0 = 2 * math.pi * f0   # rad/s
    L1 = 1.0e-3             # H (1 mH)
    L2 = 1.0e-3             # H (1 mH)
    C_res = 1.62e-9         # F (1.62 nF)
    R_esr = 10.0            # Ohm (resistencia parasita de la bobina)
    R_limit = 150.0         # Ohm (resistencia limitadora serie TX)
    R_driver_on = 3.5       # Ohm (resistencia MOSFET driver TC4427A)
    Vcc_tx = 5.0            # V
    Vdd_rx = 3.3            # V
    
    # 2. Calculo de Ganancia
    XL0 = w0 * L2
    XC0 = 1.0 / (w0 * C_res)
    Q0 = XL0 / R_esr
    Rp_ideal = (Q0 ** 2) * R_esr  # Resistencia paralelo equivalente
    
    # Carga de entrada preamplificador MCP6002 (R3 = 10k)
    Rin_amp = 10.0e3
    Rp_eff = (Rp_ideal * Rin_amp) / (Rp_ideal + Rin_amp)
    Q_eff = Rp_eff / XL0
    
    # Ganancia activa en DC
    Av1_dc = 68.0 / 10.0
    Av2_dc = 47.0 / 10.0
    Av_active_dc = Av1_dc * Av2_dc
    
    # Correccion por GBW = 1.0 MHz
    gbw = 1.0e6
    fc1 = gbw / (1.0 + Av1_dc)
    fc2 = gbw / (1.0 + Av2_dc)
    
    Av1_ac = Av1_dc / math.sqrt(1.0 + (f0 / fc1) ** 2)
    Av2_ac = Av2_dc / math.sqrt(1.0 + (f0 / fc2) ** 2)
    Av_active_ac = Av1_ac * Av2_ac
    
    Av_total_real = Q_eff * Av_active_ac

    print("\n--- GANANCIAS ---")
    print(f"Reactancia resonante: {XL0:.2f} Ohm")
    print(f"Q no cargado (Q0):   {Q0:.2f} ({20*math.log10(Q0):.2f} dB)")
    print(f"Q cargado (Q_eff):   {Q_eff:.2f} ({20*math.log10(Q_eff):.2f} dB)")
    print(f"Ganancia Activa AC:  {Av_active_ac:.2f} V/V ({20*math.log10(Av_active_ac):.2f} dB)")
    print(f"Ganancia Global:     {Av_total_real:.2f} V/V ({20*math.log10(Av_total_real):.2f} dB)")

    # 3. Balance de Potencia
    V1_pk = (2.0 * Vcc_tx) / math.pi
    V1_rms = V1_pk / math.sqrt(2)
    Z_tx_res = R_limit + R_esr + R_driver_on
    I_tx_rms = V1_rms / Z_tx_res
    
    P_limit = (I_tx_rms ** 2) * R_limit
    P_esr = (I_tx_rms ** 2) * R_esr
    P_driver = (I_tx_rms ** 2) * R_driver_on
    Q_tx_mag = (I_tx_rms ** 2) * XL0
    P_tx_total = P_limit + P_esr + P_driver + 1.25e-3 + 15.0e-3 + 115.5e-3
    P_rx_total = 3.45e-3 + 99.0e-3

    print("\n--- BALANCE DE POTENCIA ---")
    print(f"Corriente fundamental TX: {I_tx_rms*1e3:.2f} mArms")
    print(f"Potencia reactiva TX:     {Q_tx_mag*1e3:.2f} mVAR")
    print(f"Potencia total TX:        {P_tx_total*1e3:.2f} mW")
    print(f"Potencia total RX:        {P_rx_total*1e3:.2f} mW")

if __name__ == "__main__":
    main()
\end{lstlisting}

\subsection{Curvas Electromagnéticas y Respuestas Frecuenciales}

El script \texttt{generar\_graficos.py} computa la distribución espacial de inducción magnética axial $B_z(z)$ según Biot-Savart, contrasta la aproximación dipolar cúbica ($1/z^3$), calcula la inductancia mutua $M(z)$ y el factor de acoplamiento $k(z)$, y evalúa la respuesta en frecuencia de ambas redes LC.

\begin{lstlisting}[style=custompy, caption={Cálculo de campo magnético y acoplamiento (\texttt{informe\_ieee/datos/generar\_graficos.py})}]
#!/usr/bin/env python3
import os
import numpy as np
import matplotlib.pyplot as plt

MU_0 = 4 * np.pi * 1e-7      # H/m
f0 = 125.04e3               # Hz
omega0 = 2 * np.pi * f0     # rad/s
L1 = 1.0e-3; L2 = 1.0e-3    # 1 mH
C_res = 1.62e-9             # 1.62 nF
R_limit = 150.0; R_esr = 10.0
I_tx_peak = 5.0 / (R_limit + R_esr) # ~31.25 mA
radius_a = 0.015            # 15 mm radio medio
N_turns = 100               # Espiras estimadas

z_cm = np.linspace(0.5, 12.0, 300)
z_m = z_cm * 1e-2

# 1. Distribucion de campo Bz(z)
Bz_exact = (MU_0 * N_turns * I_tx_peak * (radius_a**2)) / \
           (2 * ((radius_a**2 + z_m**2)**1.5)) * 1e6 # uT
mag_moment = N_turns * I_tx_peak * np.pi * (radius_a**2)
Bz_dipole = (MU_0 * mag_moment) / (2 * np.pi * (z_m**3)) * 1e6

# 2. Inductancia mutua y acoplamiento
M_z = (MU_0 * np.pi * (N_turns**2) * (radius_a**4)) / \
      (2 * ((radius_a**2 + z_m**2)**1.5))
k_z = M_z / np.sqrt(L1 * L2)

# 3. Respuesta en frecuencia de tanques resonantes
freqs = np.linspace(50e3, 200e3, 1000)
w = 2 * np.pi * freqs
Z_tx_ser = np.abs(R_limit + R_esr + 1j * (w * L1 - 1.0 / (w * C_res)))
I_tx_ac = 5.0 / Z_tx_ser

# Tanque paralelo RX
Z_rx_par = np.abs(1.0 / (1.0 / (R_esr + 1j * w * L2) + 1j * w * C_res))
V_rx_gain = Z_rx_par / (w0 * L2)
\end{lstlisting}

\subsection{Simulación Transitoria del Enlace de Datos Completo}

El script \texttt{simular\_enlace\_datos.py} realiza una simulación en el dominio del tiempo a paso discreto (\SI{50}{\nano\second}), sintetizando una trama digital Manchester a \SI{250}{\text{bps}}, generando la portadora modulada en el tanque serie TX, calculando la fuerza electromotriz inducida en RX por acoplamiento mutuo, simulando la rectificación de envolvente y resolviendo numéricamente la ecuación diferencial del filtro RC antes del comparador LM393.

\begin{lstlisting}[style=custompy, caption={Simulación dinámica temporal del enlace OOK (\texttt{informe\_ieee/datos/simular\_enlace\_datos.py})}]
#!/usr/bin/env python3
import numpy as np

def simular_enlace_completo(distancia_cm=3.0, baudrate=250):
    # Base de tiempo de simulacion
    dt = 50e-9  # 50 ns (20 MHz sample rate)
    f_carrier = 125.0e3
    T_bit = 1.0 / baudrate
    
    # 1. Generacion de datos binarios y simbolos Manchester
    bits = [1, 0, 1, 0, 1, 1, 0, 0] # Patron de prueba
    manchester_symbols = []
    for b in bits:
        if b == 0:
            manchester_symbols.extend([0, 1])
        else:
            manchester_symbols.extend([1, 0])
            
    # 2. Sintesis de portadora modulada ASK/OOK en bobina TX
    # V_tx(t) = V_pk * sin(w0 * t) cuando simbolo == 1
    
    # 3. Acoplamiento inductivo: FEM inducida en RX
    # E_ind(t) = w0 * M(z) * I_tx(t)
    
    # 4. Amplificacion activa en MCP6002 (Av_activa = 18.63 V/V)
    
    # 5. Deteccion de envolvente no lineal (BAT85 + RC)
    # Resolver dV_env/dt = (V_rect - V_env) / (R7 * C6)
    
    # 6. Conformacion con histeresis en comparador LM393
    # V_digital = 3.3V si V_env > V_th_high, 0V si V_env < V_th_low
\end{lstlisting}

\subsection{Generador de Oscilogramas Simulados de Calidad de Instrumentación}

El script \texttt{generar\_oscilogramas\_simulados.py} sintetiza los cuatro registros de simulación en el dominio del tiempo con formato de osciloscopio digital (Figuras~\ref{fig:montaje_experimental}, \ref{fig:osc_tx}, \ref{fig:osc_rx_amp} y \ref{fig:osc_demod}), modelando los transitorios de subida del tanque serie ($\tau \approx \SI{12.2}{\micro\second}$), la polarización con masa virtual a \SI{1.65}{\volt}, la limitación de ganancia activa por GBW en el MCP6002 a \SI{125}{\kilo\hertz} y la descarga intersimbólica del filtro $R_7 C_6$.

\begin{lstlisting}[style=custompy, caption={Generador de oscilogramas simulados (\texttt{informe\_ieee/datos/generar\_oscilogramas\_simulados.py})}]
#!/usr/bin/env python3
import numpy as np
import matplotlib.pyplot as plt

# Parametros de portadora y tasa de simbolos
F0 = 125040.0; W0 = 2.0 * np.pi * F0
BAUD = 250.0; T_HALF = 0.5 / BAUD # 2.0 ms
FS = 1.25e6; DT = 1.0 / FS

def sintetizar_enlace(n_bits=8, z_cm=7.0):
    # Generacion de secuencia Manchester y moduladora GPIO18
    # Transitorios de tanque serie: tau = 2*Q / w0 ~ 12.2 us
    # Excitacion cuasiestatica V_coil = 3.18 Vpk * env * sin(w0 * t)
    # Acoplamiento mutuo y tension inducida V_tank_pp ~ 30.2 mVpp (z = 7 cm)
    # Preamplificacion MCP6002 con polo GBW (Av = 18.63 V/V sobre 1.65 V)
    # Diodo Schottky BAT85 (Vf = 0.18 V) y filtro RC (tau_env = 1.0 ms)
    # Comparador LM393 con histeresis (Vth = 0.30 V +- 25 mV)
    pass
\end{lstlisting}
